\documentclass{ctexart}
\usepackage{enumerate}
\usepackage[colorlinks,linkcolor=blue]{hyperref}

\usepackage{ulem} 
\usepackage{caption}
\usepackage{graphicx, subfig}



\title{霓虹歌姬企画}
\author{南方科技大学日语社}
\begin{document}
\maketitle
\tableofcontents

\section{前言}
前几天日语社的同学们一起去会展中心的伊势海老这个店吃日料，本来想着要不然跟服务员光说日语吧，看看会发生什么事。但是到了店里却觉得店员应该还是不可能懂日语的，遂作罢。我觉得如果想要日语炸店的话，还是要去些人均1000以上的店。虽说店员们不会说日语，可是店里依然飘荡着日语歌。几首歌听下来，觉得还挺好听的，就是我对这些歌曲和歌手都没有什么了解。同学们指出这是某某有名的歌手唱的有名的歌，我才猛然发现我虽然平素也是主要听日本的歌曲，但是对于ACG之外的日本歌曲其实并没有太多的了解。之前社里一起看红白歌会，回忆起廖同学讲着讲那讲得头头是道，加上我有点歌荒，不禁觉得是时候扩展一下我的听歌范围了。之前的我连安室奈美惠都不知道，多亏廖同学指点才算是知道了这么个人名。听歌就像看论文一样，看得多了写个综述是很好的。综述这种东西既有利于自己总结，也便于跟别人分享。遂决定组织各位写这么一个小小的综述集子，来看看日本到底有多少值得注意的歌手。每个词条大概有一张图片，再配上编写者的一些说明，最后给新人一个听歌的顺序，使得新人可以循序渐进地了解与习惯这个歌手的风格。歌手按虚拟偶像，声优，ACG歌手和非ACG歌手分类，方便读者检索。\\
“霓虹歌姬企画”这个名字是我拍脑袋定的。之所以叫歌姬是因为觉得这样叫比较萌。总之就是这样一个集子，感谢各位的贡献。——张子健


\section{虚拟偶像}
\subsection{$\mu$'s}
\begin{figure}[h]
\centering
 \includegraphics[width=1.0\textwidth]{lovelive.jpg}
 \caption{LoveLive!}
\end{figure}
$\mu$'s是LoveLive！企划里校园偶像组合的名字。“为了拯救濒临废校的学校，九位少女决定组成校园偶像”。我觉得校园偶像这个设定还是很好玩的。我刚入学的时候曾经想要为了\sout{跟妹子玩耍}宣传南科大，在学校里组织一个校园偶像团体。可是实在是找不到合适的妹子，遂作罢。\\
与同一类型的偶像大师相比，虽说LoveLive!的动画做得显然没有偶像大师765要好，但是歌还是很好听的。我个人认为LoveLive!的音乐更胜一筹。其音乐人藤泽庆昌也是少女歌剧和宝石之国的音乐人。\\
LoveLive!的乐趣还在于丰富的真人生放送节目。LoveLive!实际上是一个2.5次元企划，因为虚拟偶像背后的声优本身在企划中也有很多的亮相机会。尤其是当动画本身做得不好的时候，真人节目就变得尤为重要。各位还可以在听了一定量的歌曲之后去找$\mu$'s的演唱会看看。粉丝们的热情与小姐姐们卖力的表演可以感动任何一个有偶像之魂的人。（小姐姐一词据说就是先从LoveLive!的粉丝圈子里面流行起来的）
\subsubsection*{歌曲推荐}
实际上$\mu$'s的曲风还是比较多变的。在这里列出一些笔者喜欢的歌。其中soldier game每次去KTV都想唱...
\begin{enumerate}
\item Angelic Angel
\item No brand girls
\item soldier game
\item 春情ロマンティック
\item Music S.T.A.R.T!!
\end{enumerate}




\end{document}